% Reciprocal proofs for the disjoint marked-tree continuation fragment.
\begin{annexitem}{annex:p12-klop1987-mt-embedding}{thm:p12-klop1987-mt-embedding}{Tree classes, source embedding, and the stronger planar comparison}
\begin{proof}
Changing a sibling-order representation composes a node map with
rooted-tree isomorphisms on its domain and codomain. Such compositions
preserve injectivity, ancestry equivalence and label inequalities, so
the definition is independent of representatives. Identity gives
reflexivity. Composing two embeddings gives injectivity, ancestry
equivalence in both directions, and the transitive label inequality,
proving transitivity.

If embeddings exist both ways between finite trees, their node counts
are equal. Each node map is therefore a bijection and an isomorphism
of the ancestry orders. The sum of the natural labels cannot decrease
under either bijection. The two inequalities make those sums equal;
every individual label slack in either map is nonnegative, so every
slack is zero. An ancestry-order isomorphism preserves the root and
the immediate-parent relation, and now also all labels. Thus the
unordered tree classes are identical. This proves antisymmetry.

For the planar comparison, a diving step maps the source root into
a proper target subtree; coupling maps the root to the target root
and the child roots into separate immediate branches, in increasing
order. Induction gives an injective node map preserving and reflecting
ancestry, preserving left-to-right branch order and least common
ancestors. Conversely such a map either dives to the image-root
subtree or couples at that root, with the child images in distinct
immediate branches; restriction to those branches gives the recursive
clauses. This establishes its map characterization. Identity and
composition of these maps prove its preorder laws. Forgetting the
extra planar and meet conditions gives exactly a witnessing map for
the weaker source relation. The actual Figure~10 map violates the
meet condition while satisfying the source conditions, so their
strengths are genuinely different.
\end{proof}
\end{annexitem}

\begin{annexitem}{annex:p12-klop1987-mt-words}{thm:p12-klop1987-mt-words}{Infinite chains, products and the full minimal-bad-word argument}
\begin{proof}
Let an infinite sequence in a wqo be given. Some index has infinitely
many later members above it: otherwise, recursively choosing outside
the finitely many later successors of each already chosen index would
give an infinite bad subsequence. Restrict to that infinite successor
set and repeat the argument, which applies to every subsequence.
This constructs an infinite nondecreasing subsequence. In a product
sequence, extract such a subsequence for the first coordinate; the
wqo property of the second then supplies a comparable pair on it.
This proves the product assertion.

For words, suppose a bad sequence exists. Choose a bad sequence
\(w_0,w_1,\ldots\) with each word of least possible length among
words admitting a bad continuation after the already chosen prefix.
The countable coding breaks ties and makes each choice well-defined.
No word is empty, since an empty word embeds in every later word.
Write \(w_i=u_i a_i\), separating its last letter. From the letters
extract an infinite nondecreasing subsequence
\(a_{i_0}\leq a_{i_1}\leq\cdots\).
Consider
\(w_0,\ldots,w_{i_0-1},u_{i_0},u_{i_1},\ldots\).
It is bad. An earlier \(w_j\) embedding into a later \(u_{i_k}\)
would embed into its original \(w_{i_k}\), contradicting the
original bad sequence. If \(u_{i_k}\) embedded into \(u_{i_l}\),
append the last-letter comparison to obtain
\(w_{i_k}=u_{i_k}a_{i_k}\leq^*u_{i_l}a_{i_l}=w_{i_l}\),
the same contradiction. Thus this new sequence is a bad continuation
after the old prefix, but starts there with the shorter word
\(u_{i_0}\). This contradicts minimality. Finite words are wqo.

For \(\mathbb N\), a bad sequence would have each later number
strictly less than every earlier one, in particular be a strictly
decreasing infinite sequence of naturals. That is impossible. Thus
natural labels satisfy the wqo premise, without being finitely bounded.
\end{proof}
\end{annexitem}

\begin{annexitem}{annex:p12-klop1987-mt-kruskal}{thm:p12-klop1987-mt-kruskal}{Minimal bad trees and the child-set wqo}
\begin{proof}
Suppose there is a bad sequence of finite planar labelled trees.
Choose \(t_0,t_1,\ldots\) minimally by number of nodes at each
position among bad continuations of its chosen prefix. Finite trees
over natural labels are countably coded by finite lists of labels
and child counts, so a fixed enumeration breaks ties.
Let \(C\) be the set of all immediate child trees of the \(t_i\).
We first prove \(C\) wqo under the same planar comparison.

If it were not, choose a bad sequence of children. Assign to each
child its least parent index \(\kappa\) in the original sequence.
The children of any finite prefix form a finite set; no value repeats
in a bad sequence, since the comparison is reflexive. Its parent
indices are therefore unbounded. Choose a member whose assigned
index \(k\) is the least index occurring anywhere in this bad child
sequence, and discard its preceding finite prefix. Replace \(t_k\)
and the later original trees by this child sequence, retaining
\(t_0,\ldots,t_{k-1}\). The result is bad: an earlier \(t_i\)
embedding into one of these children would, by diving, embed into
its parent \(t_{\kappa}\) with \(\kappa\geq k>i\), contrary
to the original badness. The child suffix itself was bad.
Its first member is a proper child of \(t_k\), hence has fewer
nodes. This contradicts the minimal choice at position \(k\).
Thus \(C\) is wqo.

Each \(t_i\) is a pair of its natural root label and its finite
ordered child list in \(C^*\). The word lemma gives wqo of
\(C^*\), and the product lemma with \(\mathbb N\) gives indices
\(i<j\) with a nondecreasing root label and an increasing-subsequence
embedding of the child lists. Coupling their already existing child
embeddings gives \(t_i\preccurlyeq_{\rm pl}t_j\), contradicting
badness. This proves the planar labelled-tree theorem.

For an arbitrary infinite sequence of unordered source trees, choose
one planar representation of each. The planar theorem gives a pair
with its stronger embedding. Forgetting sibling order and the extra
meet condition gives a source embedding by the preceding lemma.
Hence the source relation is wqo. No finite bound on labels, arities
or tree depths was introduced in either argument.
\end{proof}
\end{annexitem}

\begin{annexitem}{annex:p12-klop1987-mt-embedding-reduction}{thm:p12-klop1987-mt-embedding-reduction}{The exact source map gives a marker construction}
\begin{proof}
Fix a witnessing node map. Use put then select along the path from
the target root to the image of the source root. This gives an
unmarked target subtree whose root is the source root's image.
Prove recursively that its subtree reduces to the corresponding
source subtree. Let the target root have label \(m\), and the
desired source root have label \(n\leq m\).

If \(n<m\), put a mark on that target root and copy it below a
new root \(n\), using as many marked copies as the source root
has children. Each child image is a proper descendant of the old
root, so select along its path within the assigned marked copy;
after the first select the remaining path uses put/select pairs.
This yields the image-root subtree for that source child. Apply the
recursive construction there. With no source children use zero
copies, giving the leaf \(n\).

If \(n=m\), partition the source child roots by the immediate target
branch containing their images. Source siblings have incomparable
images, by ancestry reflection. In a group of two or more, every
image root is strictly below its immediate target child: if one were
that child root, it would be an ancestor of the others. Put at the
parent and push into that target child, making one marked copy per
group member. Each marked copy can therefore select at least once
to its assigned image root and then undergo the recursive construction.
A singleton group whose image is properly below the immediate child
uses the same construction with one copy. If its image is the
immediate child itself, leave that child in place and recurse inside
it without a parent operation. A target branch with no source image
is deleted by parent put and zero-copy pushing.

The operations for different branches preserve the unmarked parent
and all other branches. They therefore concatenate in its context
and terminate after the finitely many recursive source nodes have
been processed. Every selected child root is unmarked when the
star operation is applied, and every surrounding context used in
this construction is unmarked. Thus it uses even the literal source
context closure, not just the enlarged auxiliary game.
All desired root labels and child subtrees are now exact, up to sibling
permutation. This proves \(t\to^*s\). If the endpoint classes
are different, the constructed sequence cannot have zero length.
The identity case makes no positive self-reduction claim.
\end{proof}
\end{annexitem}

\begin{annexitem}{annex:p12-klop1987-mt-tracing}{thm:p12-klop1987-mt-tracing}{All four node-origin cases and the finite permutation lift}
\begin{proof}
Give every input node a separate origin. Putting retains all origins.
Selecting retains only the chosen child's subtree and its origins.
Copying duplicates the whole marked tree with its origins; its new
lower root has the old root as origin as well, with strictly smaller
numerical label. Pushing retains the parent origin and unchanged
siblings, and copies the chosen child's node origins to every new
marked copy, or deletes that child's origins when there are zero
copies. Context nodes retain their origins. Thus each output node
has exactly one input ancestor, although one input node can have
several descendants. Composition proves unique ancestors over a
finite sequence. Label values are preserved on all transitions except
the strictly smaller new root of copying, so remain nonincreasing.

For a single step and a chosen surviving node, its rooted subtree is
unchanged, undergoes that local marker step in context, or, in a copy,
is an unchanged copied subtree. A newly copied lower root uses the
original root step itself. In pushing, a chosen child's root undergoes
put when it becomes starred; its deeper subtrees are unchanged copies.
These cases give zero or one projected step from the ancestor subtree
to the chosen descendant subtree. Composing along its unique ancestry
gives a subtree reduction of at most the original length.

Finally a sibling swap before a marker operation can be postponed.
For put, retain the swapped arguments and swap them afterwards; for
select, choose the same child occurrence in its changed position;
for copy, carry the swap into every copy; for push, choose the same
child and permute the resulting child block and unchanged siblings.
Swaps inside that child are carried into each of its copies. The
same verification applies in a surrounding context. Thus a finite
sequence of swaps followed by a marker step can be replaced by one
marker step followed by swaps. Repeating over a finite tree reduction
postpones all swaps until its endpoint and keeps the number of marker
steps. A positive cycle of unordered classes consequently lifts to
a positive planar marker sequence ending in a permutation of its
unmarked planar input. This justifies occurrence tracing without
assuming a planar order in the original source universe.
\end{proof}
\end{annexitem}

\begin{annexitem}{annex:p12-klop1987-mt-acyclic}{thm:p12-klop1987-mt-acyclic}{The minimal cycle, including deletion of an original child or a copied child}
\begin{proof}
Suppose an unmarked tree class is cyclic in the all-context game.
Choose one with the least number of nodes, and a positive cycle
from it with the least number of marker steps. Lift it to a planar
sequence \(t=n(\vec u)\rightsquigarrow^+t'\simeq t\) as above.
Its final root must have the initial root as ancestor. Otherwise its
unique ancestor is a proper node \(p\) of \(t\); subtree projection
gives \(t|_p\rightsquigarrow^*t'\). The corresponding proper
subtree of \(t'\) is a permutation of \(t|_p\). Selecting along
its proper path gives a positive reduction from \(t'\) to it.
Concatenation makes a smaller unmarked tree cyclic, a contradiction.

There is no select or copy at the root. Consider the first such step.
A first select removes the sole descendant of the original root,
contradicting its final ancestry. A first copy replaces the current
marked root \(n^*(\vec v)=a\) by a lower root \(m<n\) and
copies of \(a\). The final root label \(n\) cannot descend
from the lower \(m\), because labels never increase on a lineage.
It descends from one of the copied roots. Project the remaining
sequence to that copy: this gives \(a\rightsquigarrow^*t'\)
with no more steps than the old tail. The prefix to \(a\) is
nonempty, because the root had to be marked. Omitting the first
copy produces a shorter positive cycle on the same input, contradicting
the chosen step minimum.

Root operations are now only put and push. Each root put is paired
with the subsequent root push that removes its mark; steps between
them occur strictly below it. If there were no root push, all steps
would be below an unmarked root. They give reductions on its child
trees whose final multiset equals its initial multiset. Follow a cycle
of the finite permutation matching those child endpoints to initial
children. A child with a positive reduction lies on such a permutation
cycle; concatenating its child reductions, with their corresponding
representatives, gives a positive cycle on a proper unmarked subtree.
This again contradicts the node minimum. Thus some root push occurs.

A root push with one copy can be removed from a minimal cycle.
Omit its paired root put, perform the intervening below-root steps
with the root left unmarked, and put directly on the chosen child's
root at the former push point. All contexts are allowed in this
auxiliary game. The chosen child was unmarked at that point, so the
replacement is legal and reaches exactly the same term. Two root
steps have become one child step, contradicting the step minimum.

Consider a root push with zero copies, deleting a child. Track that
child's most recent root-level birth. Until the deletion it has had
no further root push affecting it; such a push would be a new birth.
If it was an original child, remove that child from the initial tree,
omit its internal history and the put/push pair deleting it, and
retain every operation elsewhere. This reaches the same post-deletion
term and then follows the original tail to \(t'\simeq t\).
Delete there the matching original child by put and zero-copy push.
The result is the pruned initial class. These last two steps make
the constructed cycle positive, while its initial tree has fewer
nodes, contradicting the node minimum.

If the deleted child was born as one of \(k\) marked copies of
an earlier root push, make only \(k-1\) copies at that birth.
Omit the missing copy's whole internal history and the put/push pair
that would later delete it. Other copies and siblings have identical
histories after reindexing positions. This reaches the same
post-deletion term and final class, but with fewer steps. The altered
birth push remains a legal step, including when its copy count is
zero. This construction does not assume that the erased copy had
no internal computation or was unmarked at birth. It contradicts
the step minimum. Thus zero-copy root pushes are excluded as well.

Every root push therefore makes at least two copies and strictly
increases the number of immediate root children. At least one occurs;
put and below-root operations do not change that number. The final
unmarked tree cannot be a permutation of the initial one. This final
contradiction proves acyclicity of the auxiliary all-context game,
and hence of its source subrelation.
\end{proof}
\end{annexitem}

\begin{annexitem}{annex:p12-klop1987-mt-well-founded}{thm:p12-klop1987-mt-well-founded}{The established wqo barrier and the strict source order}
\begin{proof}
Transitivity follows by concatenating positive finite reductions whose
intermediate endpoint is unmarked. Acyclicity gives irreflexivity,
so this is a strict partial order. If an infinite descending sequence
\(t_0\to^+t_1\to^+\cdots\) of unmarked trees existed, its wqo
would give \(i<j\) with \(t_i\sqsubseteq t_j\).
The embedding construction gives \(t_j\to^*t_i\).
Concatenate it with the positive reduction from \(t_i\) to
\(t_j\). The result is a positive cycle on the unmarked class
\(t_i\), contrary to the proved acyclicity. This also covers
\(t_i=t_j\): the backward part can be empty, but the forward
part is positive. Thus no infinite descending sequence exists.
\end{proof}
\end{annexitem}

\begin{annexitem}{annex:p12-klop1987-mt-boolean}{thm:p12-klop1987-mt-boolean}{Positive marker constructions for every Table~9 rule and every instance}
\begin{proof}
Encode each atomic formula or variable as the unmarked leaf 0, and
use labels 1,2,3 for disjunction, conjunction and negation. Injectivity
of this encoding is not required: we show every actual rewrite maps
to a positive marker reduction, and acyclicity then prevents a step
with identical encoded endpoints.

For double negation, select twice:
\(3(3(t))\to3^*(3(t))\to3(t)\to3^*(t)\to t\).
For \(3(1(t,s))\), put at 3, copy twice below the lower root 2,
then in each marked 3-copy push once into its 1-child. Select \(t\)
in the first resulting \(1^*\), and \(s\) in the second.
The result is \(2(3(t),3(s))\).
For \(3(2(t,s))\), use the same construction with the new lower
root 1 and marked 2-children; it gives \(1(3(t),3(s))\).

For \(2(t,1(s,r))\), put at 2 and copy twice below the lower
root 1. In each marked 2-copy push once into its 1-child, retaining
the unchanged \(t\)-child. Select \(s\) in the first \(1^*\)
and \(r\) in the second. This gives
\(1(2(t,s),2(t,r))\). For \(2(1(s,r),t)\), choose the other
displayed child in this same construction, giving
\(1(2(s,t),2(r,t))\); these choices are valid on unordered tree
classes while the source formula syntax keeps its two ordered rules.
Every construction includes a put and a copy or select, so is positive.

The placeholders \(t,s,r\) may be arbitrary encoded finite terms.
Thus the constructions apply to all substitution instances, not only
ground examples. They also remain valid in any encoded original
context, whose nodes are unmarked. An infinite original reduction
would therefore give an infinite chain of positive reductions between
unmarked trees, contrary to the proved well-foundedness. This proves
termination of all five original rules.
\end{proof}
\end{annexitem}

\begin{annexitem}{annex:p12-klop1987-mt-indexed}{thm:p12-klop1987-mt-indexed}{The precise indexed simulation and unchanged ancestor annotations}
\begin{proof}
Count every operator or variable node once. For finite terms \(s,r\),
let \(n=|s|+|r|\geq2\). Annotate each binary node recursively
using the size of its two original arguments, with labels
\(g_n=2n,h_n=2n-1\) and unary \(f=0\); variables can use
leaf 0. The annotated tree is unmarked.

For the first rule, its annotated instance is
\(g_n(t,u)\to h_n(t,u)\). Put at \(g_n\), copy twice under
the lower \(h_n\), and select \(t,u\) from the two marked
copies. This is a positive marker reduction.
For the second rule, the annotated instance is
\(h_{n+1}(f(t),u)\to f(g_n(t,u))\).
Put at \(h_{n+1}\), copy once under \(f\), then copy the
remaining marked \(h_{n+1}\) twice under \(g_n\), which is
strictly lower. In the first copy select \(f(t)\), then use
put/select there to reach \(t\); in the second select \(u\).
The final tree is \(f(g_n(t,u))\), as required.

The first original rule preserves its node count. The second has
\(2+|s|+|r|\) nodes on both sides. Therefore an enclosing binary
node sees unchanged sizes of its arguments, so its annotation index
does not change. Other subtrees keep their annotations as well.
Both simulations consequently apply inside arbitrary original
contexts, using only unmarked contexts in the marker game. An
infinite original reduction gives an infinite positive chain of
unmarked annotated trees, excluded by the main theorem. The
natural-labelled theorem covers all indices uniformly. Node-count
preservation also bounds them within any one original reduction;
no one finite bound is imposed on all inputs.
\end{proof}
\end{annexitem}

\begin{annexitem}{annex:p12-klop1987-mt-obstruction}{thm:p12-klop1987-mt-obstruction}{A reverse-embedding cycle contradiction}
\begin{proof}
If \(l\sqsubseteq r\), embedding simulation gives \(r\to^*l\).
A descending simulation \(l\to^+r\) would concatenate with it
to form a forbidden positive unmarked cycle. In the source example,
map the left f-root to the internal right-hand f, its s-child to the
s below that f's p-child, and the variable-instance subtree to its
identical copy below that s. This is an injective ancestry-order map
with equal labels on mapped operators. The obstruction is independent
of the weights of g and p. It does not exclude another termination
method, only this particular embedding-containing strict order.
\end{proof}
\end{annexitem}

